We presented a unified study of fault attacks against 
post-quantum signature schemes based on seed-tree compression. 
Through the GZKST abstraction, we identified a 
common seed-hiding invariant underlying schemes such as 
LESS-v2, CROSS, and MEDS, and showed that previously 
proposed attacks can be viewed as violations of this 
invariant.

Building on this framework, we derived generic key-recovery 
procedures and analyzed their query complexity. We also 
demonstrated practical clock-glitch attacks against the 
MEDS reference implementation on an ARM Cortex-M4 
microcontroller, revealing multiple fault surfaces that 
enable hidden-seed disclosure and secret-key recovery. 
Finally, we proposed and evaluated several countermeasures, 
showing that effective protections can be achieved with 
small overhead. Overall, our results highlight the 
importance of securing seed-publication mechanisms 
against implementation-level fault attacks.
